\documentclass[12pt]{article}
\usepackage[utf8]{inputenc}
\usepackage[T1]{fontenc}
\usepackage{hyperref}
\title{Seru Language Resources for K-12 Curriculum Development}
\author{Compiled by Claude Sonnet 5.5 (claude-sonnet-5-5) for Andr\'as Kornai}
\date{}
\begin{document}
\maketitle

\section{Introduction}
Seru (Glottocode seru1247) appears to be an extinct language of Sarawak, Malaysia, said to have died out in the 1950s. Wikipedia describes it as unattested and unwritten, and places it tentatively in the Punan group of Austronesian (Smith 2017). Glottolog marks it as insufficiently attested or not a distinct language. The ISO 639-3 code szd was reportedly retired and merged into Ukit (umi) in a 2022 change request adopted in 2023, though Glottolog still lists the code. These points come from search summaries and should be checked at the ISO 639-3 registry. Seru should not be confused with Serui-Laut (Indonesia), Serua, or Sera (Papua New Guinea).

\section{Primary Educational Resources}
None. No primer, textbook, wordlist or school programme was found.

\section{Linguistic and Reference Materials}
No grammar, dictionary or wordlist for Seru was found. Reference pages: \url{https://en.wikipedia.org/wiki/Seru_language} and the Glottolog entry \url{https://glottolog.org/resource/languoid/id/seru1247.bigmap.html}.

\section{Religious and Community Resources}
None found.

\section{Digital and Audio Resources}
None found.

\section{Curriculum Development Context}
A heritage or revitalization curriculum for Seru is not feasible from the evidence found. The language is unattested, so there is nothing to teach in it. Any heritage programme would have to be framed around the history of the people and the neighbouring Punan and Ukit communities, with expert input. I do not recommend inventing language content.
\end{document}